% No preamble. Prefix pvs-. Actual flat physical spectral estimates.
\section{Flat physical heat densities and the remaining curvature interface}
\label{pvs-section}

Fix dimension four, a color algebra of dimension \(n_{\mathfrak g}\),
a primitive odd ratio \(\ell\), a composite \(K=\ell^k\), and any
preceding homogeneous depth \(M\).
Coordinates are current-link angles with counting inner products.
On a cubic current torus of side \(KN\), there are \(N^4\) final
\(K\)-blocks.
All statements about varying backgrounds below are explicitly
separated from the flat spectral theorem.

\subsection{Actual flat inputs}
\label{pvs-input-section}

Let \(P_f,P_c\) be the fine and coarse physical projections at
the flat field, and write
\begin{equation}
 \begin{split}
 L&=q_M,\qquad A=P_cC_K|_{\operatorname{ran}P_f},\\
 R&=L|_{\operatorname{ran}P_f}-A^Tq_{MK}A,\qquad
 Q=R+aA^TA .
 \end{split}
 \label{pvs-operators}
\end{equation}
These are the actual residual physical operators of the finite
quotient identity. The following previously proved facts are the
only model-specific inputs:
\begin{equation}
 \begin{gathered}
 L\ge d^Td,\qquad 0\le q_{MK},\qquad \|q_{MK}\|\le R_q,\\
 R\ge0,\qquad Q\ge c_aK^{-2}I,\qquad c_a>0 .
 \end{gathered}
 \label{pvs-inputs}
\end{equation}
The constants are uniform in \(M,K,N\).
For example, the physical alias proof gives
\[
 c_a=\min\{B_0^{-1},a/(A_0+B_0R_q)\},\qquad
 A_0=2c_\omega^{-2},\quad
 B_0=(1+2c_\omega^{-2})/4,\quad c_\omega=(2/\pi)^5.
\]
Flat operators are invariant under final-block translations.

The globally based scalar gradient has the same range as the
ordinary periodic gradient: subtracting the value of a scalar at
the base does not change its gradient.
Thus \(\operatorname{ran}P_f\) is the usual transverse Bloch
space, including all harmonic one-forms at momentum zero.
No estimate on an inverse globally based scalar operator is used.

Choose a fixed \(a\ge R_q\). Then
\begin{equation}
 Q=L|_{\rm phys}+A^T(aI-q_{MK})A\ge\Delta_{\rm phys},
 \qquad
 Q\ge\tfrac12\Delta_{\rm phys}+\frac{c_a}{2K^2}I .
 \label{pvs-combined-order}
\end{equation}
Here \(\Delta_{\rm phys}\) is the Maxwell form on the transverse
space, where it is the scalar lattice Laplacian.

\begin{theorem}[Uniform flat physical heat and inverse moments]
\label{pvs-main}
Put \(\widehat Q=K^2Q\), and define the heat trace per final block by
\[
 \widehat H_{M,K,N}(u)
   =N^{-4}\operatorname{Tr}_{\rm phys}e^{-u\widehat Q}.
\]
There are \(C,c>0\), independent of \(M,K,N\), such that
\begin{equation}
 \widehat H_{M,K,N}(u)
 \le C\left\{\frac{K^4}{(1+K^2u)^2}+N^{-4}\right\}e^{-cu},
 \qquad u>0.
 \label{pvs-heat}
\end{equation}
In the infinite flat structural system the corresponding block
trace satisfies the same bound with \(N^{-4}\) omitted.
For every \(p>0\),
\begin{equation}
 \operatorname{Tr}_{\rm block}\widehat Q^{-p}
 \le C_p
 \begin{cases}
 K^{4-2p}+N^{-4},&0<p<2,\\
 1+\log K+N^{-4},&p=2,\\
 1+N^{-4},&p>2,
 \end{cases}
 \label{pvs-moments}
\end{equation}
again omitting \(N^{-4}\) in infinite volume.
For each fixed real \(r\) and \(U\ge1\),
\begin{equation}
 \int_U^\infty u^{r-1}\widehat H_{M,K,N}(u)\,du
 \le C_r(1+N^{-4})e^{-cU/2}.
 \label{pvs-tail}
\end{equation}
These are unperturbed spectral statements. They do not assert a
bound on a background derivative of a heat trace.
\end{theorem}

\begin{proof}
For self-adjoint finite matrices \(X\ge Y\), the min--max theorem
gives ordered eigenvalues \(\lambda_j(X)\ge\lambda_j(Y)\);
therefore \(\operatorname{Tr}e^{-sX}\le
\operatorname{Tr}e^{-sY}\).
No operator monotonicity of the exponential is asserted.

At momentum \(p\in(2\pi/n)\mathbb Z_n^4\), \(n=KN\), put
\(P(p)=4\sum_\mu\sin^2(p_\mu/2)\).
The physical Maxwell multiplicity is \(3n_{\mathfrak g}\)
away from zero and \(4n_{\mathfrak g}\) at zero. Hence
\begin{equation}
 \operatorname{Tr}_{\rm phys}e^{-sQ}
 \le 4n_{\mathfrak g}e^{-c_as/(2K^2)}
                       \sum_p e^{-sP(p)/2}.
 \label{pvs-minmax}
\end{equation}
For the one-dimensional scalar heat average
\[
 p_n(t)=\frac1n\sum_{m\ {\rm mod}\ n}
                      e^{-4t\sin^2(\pi m/n)},
\]
use representatives \(|m|\le n/2\) and
\(\sin(\pi|m|/n)\ge2|m|/n\).
Gaussian summation for \(t\ge1\), and \(p_n(t)\le1\) for
\(t\le1\), give
\begin{equation}
 p_n(t)\le n^{-1}+C(1+t)^{-1/2},\qquad
 p_n(t)^4\le C\{n^{-4}+(1+t)^{-2}\}.
 \label{pvs-scalar-heat}
\end{equation}
Since \(\sum_p e^{-tP(p)}=n^4p_n(t)^4\),
division of \eqref{pvs-minmax} by \(N^4\) yields
\[
 N^{-4}\operatorname{Tr}e^{-sQ}
 \le C\{K^4(1+s)^{-2}+N^{-4}\}e^{-cs/K^2}.
\]
Set \(s=K^2u\) to obtain \eqref{pvs-heat}.

The normalized total trace equals the diagonal trace in any one
final block, by final-block translation invariance.
This equality does not turn the min--max argument into a
pointwise operator comparison.
In infinite volume, apply the same min--max argument in each
finite-dimensional coarse Bloch fiber and integrate.
Folding the Maxwell aliases gives \(K^4\) times the fine scalar
heat density, proving the infinite statement directly.

Positive spectral calculus gives
\[
 \operatorname{Tr}_{\rm block}\widehat Q^{-p}
  =\frac1{\Gamma(p)}\int_0^\infty
      u^{p-1}\operatorname{Tr}_{\rm block}e^{-u\widehat Q}\,du .
\]
Split at \(K^{-2}\) and \(1\) to obtain
\eqref{pvs-moments}.
For \(p>2\), one can instead use
\(K^4/(1+K^2u)^2\le u^{-2}\) directly.
The same inequality and exponential decay prove
\eqref{pvs-tail}.
\end{proof}

The finite \(N^{-4}\) term is necessary: it retains the finite
number of physical harmonic modes.
The local inverse moment means
\(\operatorname{Tr}(P_{\rm block}\widehat Q^{-p}P_{\rm block})\).
It is not a claim of global trace class in infinite volume, nor
an identification with a power of the compressed inverse.
The logarithm in the \(p=2\) bound is an upper bound; no matching
actual asymptotic is claimed.

\subsection{An optional finite on-shell replacement}
\label{pvs-onshell-section}

On a finite real on-shell family with
\(LD_f=0\) and \(qD_c=0\), set
\[
 Q_{\rm on}=L|_{\rm phys}+aA^TA=Q+A^Tq_{\rm phys}A.
\]
The minimum right inverse gives \(AQ^{-1}A^T=a^{-1}I\).
Consequently the determinant lemma proves
\begin{equation}
 \frac{\det Q_{\rm on}}{\det Q}
   =\det(I+a^{-1}q_{\rm phys}).
 \label{pvs-onshell-determinant}
\end{equation}
Because \(q\) kills the coarse gauge space, the normalized
logarithm on the right is the normalized full coarse
\(\log\det(q+aI)\); only a field-independent power of \(a\)
has been removed.
The already proved coarse-factor Maxwell estimate can therefore
control this replacement wherever its Maxwell moment is
legitimately defined.

This does not apply to an arbitrary off-shell family by omission
of the multiplier curvature.
The actual general form is
\(L=s''-\lambda\cdot f''\).
An infinite affine use of separate globally based physical
determinants requires its own infrared justification.

\subsection{A precise sufficient derivative estimate}
\label{pvs-vertex-section}

The following fixed-space lemma isolates an additional estimate
which would suffice for an infrared second-variation bound.
Its hypotheses are not established here for the actual varying
physical quotient.

\begin{lemma}[Relative-form vertices]
\label{pvs-vertex-lemma}
Let
\(X(t)=X+tV_1+t^2V_2+O(t^3)\)
be a real self-adjoint family on a fixed finite-dimensional
Hilbert space, with \(X>0\). Suppose
\begin{equation}
 \|X^{-1/2}V_jX^{-1/2}\|\le L_j,\qquad j=1,2.
 \label{pvs-vertex-hyp}
\end{equation}
Then
\begin{equation}
 \left|[t^2]\operatorname{Tr}e^{-uX(t)}\right|
 \le8(L_2+L_1^2)\operatorname{Tr}e^{-uX/4},
 \qquad u>0.
 \label{pvs-vertex-bound}
\end{equation}
\end{lemma}

\begin{proof}
The single insertion has absolute value at most
\[
 uL_2\operatorname{Tr}(Xe^{-uX})
 \le(2/e)L_2\operatorname{Tr}e^{-uX/2}.
\]
In an eigenbasis of \(X\), cyclic Duhamel gives the double
insertion exactly as
\begin{equation}
 \frac u2\sum_{i,j}|(V_1)_{ij}|^2
       \frac{e^{-u\lambda_i}-e^{-u\lambda_j}}
                         {\lambda_j-\lambda_i},
 \label{pvs-divided-difference}
\end{equation}
with the continuous value \(ue^{-u\lambda_i}\) at coincident
eigenvalues.
For \(a,b>0\),
\begin{equation}
 ab\,\frac{e^{-ua}-e^{-ub}}{b-a}
 \le2\{ae^{-ua/2}+be^{-ub/2}\}.
 \label{pvs-elementary}
\end{equation}
Indeed, assume \(b\ge a\).
For \(b\ge2a\), use \(ab/(b-a)\le2a\).
For \(a\le b\le2a\), use the derivative bound
\(ue^{-ua}\), \(ab\le2a^2\), and
\((ua)e^{-ua/2}\le2/e\).

Write \(V_1=X^{1/2}BX^{1/2}\), with \(\|B\|\le L_1\).
The row and column square sums of \(B\) are at most \(L_1^2\).
Applying \eqref{pvs-elementary} to
\eqref{pvs-divided-difference} bounds it by
\[
 2uL_1^2\operatorname{Tr}(Xe^{-uX/2})
 \le(8/e)L_1^2\operatorname{Tr}e^{-uX/4}.
\]
Together with the single insertion this proves the stated
conservative constant.
\end{proof}

If \(X(t)=K^2Q_\perp(t)\) could be identified with a fixed
physical space and satisfy \eqref{pvs-vertex-hyp} uniformly,
Theorem~\ref{pvs-main} would give
\begin{equation}
 \int_1^\infty
 \left|N^{-4}[t^2]\operatorname{Tr}e^{-uX(t)}\right|
 \frac{du}{u}
 \le C(L_2+L_1^2)(1+N^{-4}).
 \label{pvs-conditional-tail}
\end{equation}
An infinite-volume version additionally requires a justified
trace-and-derivative limit.
Neither the uniform identification of varying physical spaces
nor the actual bounds \eqref{pvs-vertex-hyp} follows from the
flat gap.
Derivatives of \(P_f,P_c\), the current, and the nonlinear
constraint cannot be absorbed into an unproved constant.

\subsection{The required scalar--physical cancellation}
\label{pvs-scope-section}

The separate physical determinant has a growing response even
along constant commuting final link angles.
The actual globally based finite identity and the local Ward
limit give, at fixed \(K\),
\begin{equation}
 \lim_{N\to\infty}N^{-4}[t^2]\log\det Q_\perp(tA)
   =C_H|A|^2(K^2-1)g_4,\qquad
 g_4=\int_{[-\pi,\pi]^4}
       \frac{dp}{(2\pi)^4\,4\sum_\mu\sin^2(p_\mu/2)}>0 .
 \label{pvs-holonomy}
\end{equation}
Here the fine angles are \(tA/K\).
This is the separately established flat-holonomy calculation:
the same term in the scalar orbit-density ratio cancels
\eqref{pvs-holonomy} in the actual total correction.
It is a zero-momentum test, not a magnetic Maxwell coefficient.

Accordingly, the spectral theorem does not permit an independent
replacement of a massless physical or scalar factor by its
massive infinite reference.
A valid curvature comparison must preserve the coupled
cancellation, or supply a justified joint infrared regulator.
No actual vector logarithmic coefficient, its sign, or
Yang--Mills existence conclusion is proved by this section.

